\documentclass[manuscript,anonymous]{acmart}

\usepackage{tcolorbox}
\usepackage{subcaption}
\usepackage{enumitem}
\usepackage{multirow}
\usepackage{algorithm}
\usepackage{algpseudocode}
\usepackage{graphicx}
\usepackage{textcomp}
\usepackage{listings}
\usepackage{url}
\usepackage{xurl}
\usepackage{amsmath}

\input{solidity_highlight}

\newcommand{\tool}[1]{\texttt{#1}}

\newtcolorbox{rqbox}{colback=gray!8, colframe=gray!60, boxrule=0.5pt, arc=2pt,
  left=6pt, right=6pt, top=3pt, bottom=3pt}
\newtcolorbox{answerbox}{colback=green!4, colframe=green!45!black, boxrule=0.5pt,
  arc=2pt, left=6pt, right=6pt, top=3pt, bottom=3pt}

\begin{document}

\title{VeriPUT: Verifier-Derived Parameterized Unit Tests for Path Coverage of
Smart Contracts}
\begin{abstract}
\end{abstract}
\maketitle

\section{Introduction}\label{sec:intro}

Smart contracts hold digital assets on public blockchains, and a fault in one
can be turned directly into a significant loss~\cite{werner2022sok}. Once such
contracts are deployed, they cannot be edited in place, so every version has to
be checked before it goes on chain. Developers write unit tests for this
purpose. One common way to obtain such tests automatically is \emph{coverage-based
test generation}, which measures how much of a contract a suite reaches and
generates tests that raise that
measure~\cite{syntest_solidity,solar_sbst,SolTG-1}. A suite is also kept after
it is written and run again whenever the code changes. Optionally, it carries assertions that record how the contract behaves
now, so a change that alters this behavior makes the suite fail.

\textbf{Challenges.} In general, coverage-based test generation for smart
contracts rests on the assumption behind coverage criteria, that a suite
reaching more of a contract detects more of its faults~\cite{AmmannOffutt}.
However, reaching a fault is not the same as detecting it. In fact, coverage can
be a weak indicator of how many faults a suite detects. A suite can reach every
line of a program and still detect almost none of the faults injected into the
code it reached~\cite{Wang2025MutationGuided}. The reason is that covering a
target only requires one input that reaches it, and such an input is what these
tools deliver. A contract is checked once and then released, so a fault its
suite does not notice is deployed with the code. Specifically, we identify three
obstacles between reaching a target and detecting a fault at it. The first two
concern the test that a generator emits for a target it reaches. The third
concerns the targets for which it emits nothing, and what their absence tells a
developer before deployment.

\textbf{(1) Weak oracles.} A test is expected to carry an oracle that a fault
can violate. However, a coverage criterion states which parts of a contract a
run must reach, and nothing about what should hold once it reaches them. A
coverage-driven generator therefore has no source of truth beyond the program
under test, and the strongest assertion it can synthesize is a record of the run
it produced~\cite{Barr2015,syntest_solidity}. Such a test reports the faults
that stop the run, for instance a changed function signature that breaks
compilation, or a changed \texttt{require} condition that makes the call revert.
However, a more subtle logical fault lets the run finish and only moves a value,
such as a rounding direction or a bound shifted by one. The run still reaches
the statement that carries it, so the coverage figure is unchanged, and the
fault goes unreported unless the recorded values happen to include the one it
moved. For a smart contract this second kind is what a suite is kept for. The
values a call leaves behind are the balances and stored amounts that later
transactions read, so an edit that moves one has to be reported whether or not
it is a vulnerability.

\textbf{(2) Points instead of ranges.} A test is expected to run on the inputs
at which a fault becomes visible. However, a coverage criterion is satisfied by
one input per target, and it says nothing about which input that should be. A
generator that maximizes coverage therefore has no reason to run a second input
against a target it has already reached. It also has no basis for preferring one
input over another, since it takes its targets from the structure of the program
and does not predict where a fault lies. This leads to a coverage-safety gap. In
particular, existing coverage-driven generators for smart contracts emit a
\emph{concrete replay test}, one in which every parameter is fixed to a single
value. Such a test runs the target on one of the values that reach it. The other
values are never run. If the contract behaves differently on one of them, the
test cannot see it, and the coverage figure is again unchanged. For a smart
contract, any account can call the function with a value the suite never ran.

\textbf{(3) Unreached is not unreachable.} A generator is also expected to
account for the targets it leaves uncovered, and to separate those that no input
can reach from those it failed to find. However, a coverage criterion records
only whether a target was reached. That is, for a target left uncovered, we do
not know whether any input reaches it at all. A search-based generator stops when
its budget runs out, and a symbolic one stops when solving fails, and in practice
neither outcome tells the two apart. Raising the budget does not settle the
question either, since it moves a few targets into the covered set and reports
the rest exactly as before. A search also starts from seeds and explores in an
order that changes between runs, so a target missed in one run may be reached in
the next. For a smart contract this is the part that matters. A target that was
missed by accident is behavior no test ran, and any transaction can reach it once
the contract is on chain.

\textbf{Our approach.} To address these challenges, we present \tool{VeriPUT}, a
path-coverage-driven, verifier-certified framework that generates
\emph{parameterized unit tests} (PUTs) for smart contracts. A PUT declares the
test arguments instead of fixing them, states an assumption that says which
values the test admits, and carries an assertion that must hold for every
admitted value~\cite{Tillmann2005}. The core idea has three parts. (i) We
widen how an execution path is covered, from a concrete replay test that runs
the path on one input, to a range of inputs that all follow the same path, i.e.,
parameterized inputs. (ii) We widen what the test observes at the end of the
path, from the single value one run produced, to whether each observed quantity
changes and what the quantity becomes when it changes. (iii) We use a verifier
with unbounded proving power as the oracle, which certifies that the resulting
PUT holds on the contract under test. As a result, \tool{VeriPUT} returns one test per path, and reports each path for which \tool{VeriPUT}
produced no test as proved out of reach or left undecided.

\textbf{Workflow.} \tool{VeriPUT} takes a contract and works in three stages.
The first stage instruments each entry point so that a run records the decisions
it takes. \tool{VeriPUT} then takes one path at a time and plants at each exit
an assertion stating that the recorded decisions are not those of that path. A
verifier that refutes such an assertion hands back an input that reaches the
path. Each path is therefore refuted with a counterexample, proved out of reach,
or left undecided where the verifier exhausts its limits, and \tool{VeriPUT}
reports the three separately. The second stage widens that counterexample into a
\emph{certified input region}, a range of inputs around the counterexample, and
has the verifier certify that every input in the region follows the same path.
Where no region can be certified, the region narrows to the counterexample
alone, and the test stays concrete. The third stage puts candidate assertions to
the verifier, with the region as the assumption, and writes into the test those
the verifier proves. Where none of them is proved, the test asserts the exact
values recorded in the counterexample, so every emitted test carries an
assertion the verifier has established.

\textbf{Observations.} This design is based on two observations. (i) Unlike many
other programs, a smart contract has a fixed and uniform set of behaviors to
observe and test. A test can observe the balances and stored values of the
accounts involved, the data returned, whether the call succeeded, and the log
entries produced. Assertion candidates can therefore be enumerated directly from
the contract's own state variables, return types, and events. (ii) A smart
contract commonly has many entry points, i.e., its public and external
functions, and each entry point is a short unit. This makes it suitable to
analyze the paths through each function unit on its own, rather than the paths
that run from a single entry to an exit as in many programs written in other
languages. A typical entry point checks a few conditions, updates a few state
variables, and returns. The decisions such a unit takes are few enough that the
paths through it can be enumerated one by one.

\textbf{Contributions.}
In summary, our contributions are:
\begin{itemize}
\item \textbf{Path-coverage-driven PUT synthesis.}
We present \tool{VeriPUT}, a method that generates parameterized unit tests for
smart contracts from the contract source alone, without a written specification
or an existing test suite. Compared to the concrete replay test generated by the
existing tools, a test emitted by \tool{VeriPUT} summarizes whether a path can
be taken, what the call does along it, and how the call ends, and thereby covers
an execution path in a stronger sense.

\item \textbf{Certified input regions.}
We introduce a region construction algorithm that turns the single input a
verifier returns for a path into a \emph{certified input region}, i.e., a range
of inputs proven to drive execution along that same path. Specifically, it
introduces a candidate ladder that searches the range of each argument in one
batch of verifier queries. Additionally, it introduces a practical degradation step that ranges over the
arguments it can generalize and fixes the rest, which avoids the path
combination explosion. The
resulting PUTs differ from existing coverage-driven generators for smart
contracts, which emit a single input for each path they reach.


\item \textbf{Verified assertion synthesis.}
We introduce an assertion oracle synthesis algorithm that synthesizes
path-level assertions. Specifically, for one path we observe whether the path
can be taken, whether it changes the contract's observable state, and what that
state becomes, and the algorithm asserts all three. This extends what coverage
means while staying coverage-driven, and closes the coverage-safety gap. It does
so without predicting where a fault may lie and without mutants of the contract,
which sets it apart from methods that target contract-level properties.

\item \textbf{Implementation \& Evaluation.}
Our evaluation results show that \dots

\end{itemize}

%-----------------------------------------------------
\section{Background}\label{sec:background}
%-----------------------------------------------------

%-----------------------------------------------------
\subsection{Smart Contracts and Solidity}
%-----------------------------------------------------

A smart contract is a program deployed on a blockchain, on Ethereum usually
written in Solidity~\cite{solidity_docs} and run by the Ethereum Virtual
Machine (EVM)~\cite{wood_yellowpaper_evm}. A contract declares \emph{state
variables}, whose values make up the contract \emph{state}. The state is kept
in storage between calls, so what a call does depends on its arguments and on
what earlier calls left behind. Functions read and change the state, and those
marked \texttt{public} or \texttt{external} are the \emph{entry points} of the
contract~\cite{solidityStructureOfAContract}. A contract normally has no single fixed
entry point and does not start on its own. It runs only when an outside party
sends it a \emph{transaction}, which names the entry point to call, carries its
arguments, and may transfer funds with the call. Nothing restricts how often or
in what order the entry points are used. A transaction is all or nothing. It
either finishes and keeps its writes, or it \emph{reverts} and every write it
made is undone~\cite{liu2025state}. Such contracts often hold financial
assets~\cite{werner2022sok}, and deployed code cannot be changed
afterwards~\cite{wood_yellowpaper_evm}, so a fault found later cannot simply be
patched. Therefore, a contract has to be checked before it is deployed.

The first four bytes of the call data are a \emph{function selector} that picks
the entry point, and the arguments behind it follow the \emph{application
binary interface} (ABI) of the contract~\cite{abi_spec}. They are decoded
before the function body is entered, so call data that does not fit the
declared parameter types can be rejected before the body runs. The body need
not reach its end either, as a runtime error such as a division by zero aborts
the call~\cite{solidity-panic}. A contract may also call another contract whose
address it holds. The code at that address then runs, and the caller does not
control it. The call may fail instead of returning, and the callee may call
back into the caller before the first call
returns~\cite{solidity_address_members,solidity2}.

%-----------------------------------------------------
\subsection{Unit Testing and Coverage Criteria}
%-----------------------------------------------------

A \emph{unit test} runs one unit of a program on some inputs and judges the
run. The part that judges is the \emph{test oracle}~\cite{Barr2015}. It is
often an \emph{assertion} in the test, but need not be. A test with no
assertion still fails when the run crashes, and this implicit oracle is all
that many generated tests carry~\cite{Barr2015}. A test with fixed inputs also
fixes one expected outcome, and so covers a single point of the input
space~\cite{Tillmann2005}. A \emph{parameterized unit test} turns the inputs
into parameters instead, stating an \emph{assumption} that says which values
are admitted and an assertion that must hold for all of them; the framework
supplies the concrete values, for example by
fuzzing~\cite{Tillmann2005,Claessen2000}.

A \emph{coverage criterion} measures how thoroughly a suite exercises a
program. Statement and branch coverage, for example, ask for every statement
and every branch outcome to be executed, while \emph{path coverage} asks for
every execution path to be taken~\cite{Howden1976}. Path coverage is the
strongest of the three, as it subsumes statement and branch coverage and asks
for every combination of decisions to be taken. However, paths grow
exponentially with the number of branches, and their number need not be finite
once loops appear. Some paths are \emph{infeasible} as well, in that no input
drives execution along them, and whether a given path is feasible cannot be
decided in general~\cite{AmmannOffutt}. Coverage is thus reported against the
feasible paths. A criterion says which parts of a program a suite reaches, not
whether the suite would notice a fault in them.

To run unit tests on a smart contract and collect coverage, developers use a testing
framework such as \tool{py-evm}~\cite{pyevm}, \tool{Remix}~\cite{remix},
\tool{Truffle}~\cite{truffle}, \tool{Hardhat}~\cite{hardhat}, or
\tool{Foundry}~\cite{foundry}. Such a framework deploys the contract on a local
chain, lets each test set the caller and the value sent, and supplies the
arguments of a parameterized test. An oracle written against a contract need
not check a return value. A test may require that a call revert, or watch a
balance or a state variable for an expected change.

%\emph{Mutation testing} asks whether a
%suite would notice a fault, by injecting small syntactic changes, each yielding
%a \emph{mutant} of the program, and counting how many of them make the suite
%fail~\cite{Howden1982,mutationSurveySC}. A mutant that does is killed; one that
%behaves identically to the original on every input is equivalent, and no suite
%can kill it.


%-----------------------------------------------------
\subsection{Bounded Model Checking and Counterexamples}
%-----------------------------------------------------

Formal verification checks by mathematical means whether a program satisfies a
property stated in a formal language~\cite{BaierKatoen2008}. Model checking
answers such a question by examining the states a program can
reach~\cite{Clarke1999ModelChecking}. Bounded model checking (BMC) narrows the
question to a fixed number of steps, asking whether the program can violate the
property within that many steps~\cite{Biere1999}. It unrolls the executions of
the program up to that bound and hands the question to a satisfiability solver.
A satisfiable answer means a violation exists inside the bound. An
unsatisfiable one means none does, but says nothing about longer executions, so
a property that survives is still not shown to hold in general. The solver may
also exhaust its resources and return no answer. \emph{k-induction} can lift
the bound~\cite{gadelha2017kinduction}. It proves two things. The base case is
that the property is not violated within the bound. The inductive step is that
once the property has held for a number of consecutive steps, it holds for the
next one too. Where both are proved, the result carries to every reachable
state.

The property a verifier is asked about may be written by the developer as an
assertion, or come from the verifier itself, which adds checks for arithmetic
overflow, division by zero, out-of-bounds accesses and the like~\cite{smtchecker}.
\emph{Assumptions} work the other way and narrow the inputs the verifier
considers. When the question is satisfiable, the satisfying assignment records
how the property breaks, and a verifier can return it as a
\emph{counterexample}: input values, the conditions along the path taken, an
execution trace, and the values of program variables where the property
fails~\cite{counterexample}.

\section{Motivation}\label{sec:motivation}

We use a fee-charging vault and two faults injected into it. Our goal is to show
(i)~what a coverage-driven generator returns for such a contract, (ii)~how a
concrete replay test and a parameterized unit test differ on the same path, and
(iii)~the coverage-safety gap, in that both tests cover that path while the
assurance each of them gives can differ.

\begin{lstlisting}[language=Solidity, numbers=left, numberstyle=\scriptsize,
xleftmargin=1.5em, basicstyle=\ttfamily\scriptsize, breaklines=true,
columns=fullflexible, keepspaces=true, float=tbp,
caption={A simplified fee-charging vault, with the two injected faults marked
at the line they replace.},
label={lst:motivation}]
pragma solidity ^0.8.0;
contract FeeVault {
    address public owner = msg.sender;
    address public feeReceiver = msg.sender;
    uint16 public constant feeBps = 250;    // 2.5% fee
    uint256 public constant maxFee = 1e17;  // cap 0.1 ether
    mapping(address => uint256) public deposits;
    mapping(address => uint16) public discountBps;
    event Deposited(address indexed u, uint256 amount);
    event Withdrawn(address indexed u, uint256 net, uint256 fee);

    function deposit() external payable {
        deposits[msg.sender] += msg.value;
        emit Deposited(msg.sender, msg.value);
    }

    function setDiscount(address u, uint16 bps) external {
        require(msg.sender == owner && bps <= feeBps);
        discountBps[u] = bps;
    }

    function withdraw(uint256 amount) external returns (uint256 net) {
        uint256 d = deposits[msg.sender];
        require(amount > 0 && amount <= d);
        uint256 rate = feeBps;
        if (discountBps[msg.sender] > 0) {
            rate = feeBps - discountBps[msg.sender];
        }
        // M1: / 10000 -> / 100
        // M2: amount * rate -> d * rate
        uint256 fee = amount * rate / 10000;
        if (fee > maxFee) {
            fee = maxFee;
        }
        net = amount - fee;
        deposits[msg.sender] = d - amount;
        (bool ok1, ) = feeReceiver.call{value: fee}("");
        require(ok1);
        (bool ok2, ) = msg.sender.call{value: net}("");
        require(ok2);
        emit Withdrawn(msg.sender, net, fee);
    }
}
\end{lstlisting}

Listing~\ref{lst:motivation} shows the simplified vault, which we take as the
reference program, that is, as the version whose behaviour the developer
intended. The vault holds Ether for the accounts that deposit it, charges a fee
on each withdrawal, and lets the owner grant an account a discount on that fee.
The deployer owns the contract and collects its fees (Lines~3--4). A
withdrawal is charged 250 basis points, or $2.5\%$, of the amount taken, and no
single fee exceeds 0.1 Ether (Lines~5--6). Each account has a deposit and a
discount expressed in the same unit as the fee, zero unless the owner grants one
(Lines~7--8, 17--20), and depositing credits the caller (Lines~12--15). Two
events record a deposit and a withdrawal (Lines~9--10). A withdrawal admits any
positive amount within the caller's deposit (Lines~23--24), lowers the rate by
the discount when one is held (Lines~25--28), and applies the rate, lowering the
fee to the cap when the fee exceeds it (Lines~31--34). It then computes the net,
reduces the deposit by the amount (Lines~35--36), pays the fee and then the net,
requiring each transfer to succeed (Lines~37--40), and emits the withdrawal
event (Line~41). Each conditional has two outcomes, and all four
combinations are reachable, so a successful withdrawal follows one of four
paths.

We now mutate the program, that is, we deliberately introduce an edit of a kind
a developer might make and use it to ask whether a test detects the change in
logic it causes. We introduce two such faults. Both replace Line~31, and both
are marked there in Listing~\ref{lst:motivation}. $M_1$ divides by one hundred
rather than by ten thousand, a confusion between a percentage and a basis
point, so every fee is a hundred times the intended one, or the cap where that
exceeds it. $M_2$ applies the rate to the balance rather than to the amount
withdrawn, which is what the line would have read in a version that only ever
withdrew everything, so a partial withdrawal is charged on the whole deposit.
Both edits can change which branch a given input takes, since the conditional
below them reads the fee they change. Neither edit, however, changes the text
of either condition itself.

\begin{figure*}[t]
\begin{minipage}[t]{0.46\textwidth}
\begin{lstlisting}[numbers=none, basicstyle=\ttfamily\scriptsize,
xleftmargin=0.5em, breaklines=true, columns=fullflexible, keepspaces=true,
caption={The counterexample (simplified) returned for one condition by
\tool{CC-SolBMC}~\cite{CCSolBMC2024} via \tool{SMTChecker}~\cite{smtchecker}.},
label={lst:motivation-ce}]
BMC: Assertion violation happens here.
  --> FeeVault_mod.sol:32:2
   |
32 | assert(!(discountBps[msg.sender] > 0));
   | ^^^^^^^^^^^^^^^^^^^^^^^^^^^^^^^^^^^^^
Counterexample:
  amount                  = 21239
  deposits[msg.sender]    = 21239
  discountBps[msg.sender] = 39
  ...
\end{lstlisting}

\vspace{0.4em}

\begin{lstlisting}[language=Solidity, numbers=left, numberstyle=\scriptsize,
basicstyle=\ttfamily\scriptsize, xleftmargin=1.5em, breaklines=true,
columns=fullflexible, keepspaces=true,
caption={The same counterexample rendered as a concrete-replay
\tool{Foundry}-style test unit.},
label={lst:motivation-tc}]
import "forge-std/Test.sol";
import "../src/FeeVault.sol";

contract CoverageTest is Test {
    FeeVault v;
    address alice = address(0xA11CE);

    function setUp() public { v = new FeeVault(); }
    receive() external payable {}

    function test_discountedWithdrawal() public {
        v.setDiscount(alice, 39);
        vm.deal(alice, 21239);
        vm.startPrank(alice);
        v.deposit{value: 21239}();
        v.withdraw(21239);
        vm.stopPrank();
    }
}
\end{lstlisting}
\end{minipage}
\hfill
\begin{minipage}[t]{0.50\textwidth}
\begin{lstlisting}[language=Solidity, numbers=left, numberstyle=\scriptsize,
basicstyle=\ttfamily\scriptsize, xleftmargin=1.5em, breaklines=true,
columns=fullflexible, keepspaces=true,
caption={Listing~\ref{lst:motivation-tc} widened into a parameterized unit test
over the same path.},
label={lst:motivation-put}]
import "forge-std/Test.sol";
import "../src/FeeVault.sol";

contract DiscountTest is Test {
    FeeVault v;
    address alice = address(0xA11CE);

    function setUp() public { v = new FeeVault(); }
    receive() external payable {}

    function test_discountLowersFee(uint256 dep, uint256 amount, uint16 disc) public {
        dep = bound(dep, 21239, 200000);
        amount = bound(amount, 10000, 21239);
        disc = uint16(bound(uint256(disc), 1, 100));
        v.setDiscount(alice, disc);
        vm.deal(alice, dep);
        vm.startPrank(alice);
        v.deposit{value: dep}();
        uint256 net = v.withdraw(amount);
        vm.stopPrank();
        uint256 fee = amount - net;
        assertLe(fee, amount * 249 / 10000);
        assertGe(fee, amount * 150 / 10000);
    }
}
\end{lstlisting}
\end{minipage}
\caption{What a coverage-driven generator produces for
Listing~\ref{lst:motivation} (left) and a test over the same path that reports
both faults (right). The generator targets one atomic condition at a time and
returns an assignment of values; Listing~\ref{lst:motivation-tc} is that
assignment rendered as a call by us, since the generator emits no executable
artifact. Both test contracts deploy the vault themselves and so collect its
fees, which is why each declares \texttt{receive}. The harness supplies
\texttt{vm.deal}, which credits an account with Ether;
\texttt{vm.startPrank} and \texttt{vm.stopPrank}, which make the calls between
them come from a given account; \texttt{bound}, which maps an argument into a
range; and \texttt{assertLe} and \texttt{assertGe}, which fail the test when
the ordering they state does not hold.}
\label{fig:motivation-tests}
\end{figure*}

We first ask what a coverage-driven generator produces for this contract, using
an existing one. \tool{CC-SolBMC}~\cite{CCSolBMC2024} works at the level of
atomic conditions: for every condition it finds in the source it emits a test
input that makes the condition true and one that makes it false. It obtains
each of them by planting a pair of assertions on the line above the condition,
one asserting that the condition does not hold and one asserting that its
negation does not hold. An assertion of the first kind is false exactly when
the condition is true, so asking a verifier to prove it is the same as asking
whether any input satisfies the condition. Refuting it establishes that the
true branch is reachable and returns a witness.

For the discount test of the vault (Line~26 of
Listing~\ref{lst:motivation}) the pair reads
\texttt{assert(!(discountBps[msg.sender] > 0))} for the true polarity and
\texttt{assert(!(!(discountBps[msg.sender] > 0)))} for the false one. We follow
the first. The instrumented contract is handed to
\tool{SMTChecker}~\cite{smtchecker}, the verifier built into the \tool{Solidity}
compiler, whose bounded model checking engine \tool{CC-SolBMC} drives, and the
assertion is refuted. Listing~\ref{lst:motivation-ce} shows the resulting
counterexample, the one for the true polarity, simplified to the assignments
that can serve as test inputs. It names the line and the assertion that failed
and then assigns a value to each variable the engine tracked; of those, three
determine the call. The caller has 21239 wei on deposit, holds a discount of 39
basis points, and asks to withdraw the whole balance. Those three values are
what a test needs, and we render them as a \tool{Foundry} test in
Listing~\ref{lst:motivation-tc}. The test contract deploys the vault, so it is
the owner and grants \texttt{alice} the discount of 39 basis points
(Lines~4--9, 12). It then funds \texttt{alice} with the deposit recorded in the
counterexample, deposits it, and withdraws all of it under her identity
(Lines~13--17). Note that both this test and the one beside it are written
against the reference contract, not against either fault, and that both have
been run and pass on it.

We now run this test unit against the two mutants. On $M_1$ it fails. The fault
multiplies every fee by one hundred, so at the recorded input the fee comes to
44814 wei against a deposit of 21239, the subtraction on Line~35 of
Listing~\ref{lst:motivation} underflows, and the call reverts. The test carries
no assertion, so what reported the fault is the run not completing. Notably, a
test reports a fault when the faulty program yields a different observation.
Whether a call reverted is one of the things a caller can observe, alongside
the stored state and the value returned. The test was not written with $M_1$ in
mind, and reporting a fault it was not written for is what a regression suite
is for. On $M_2$ the test passes, and no assertion added to it would have
changed that. The counterexample withdraws the whole balance, so the amount and
the deposit coincide, and charging the rate on one is charging it on the other.
The two programs agree on the fee, on the net, on the stored deposit, and on
the event. There is nothing on this run for an assertion to separate, and the rest of the
suite the generator returned for this contract fares no better. In most of its
tests the amount requested equals the deposit, and $M_2$ charges what the
reference charges. The remaining ones withdraw part of the balance, but not on
this path: one is small enough that both fees floor to zero, and two are large
enough that both are lowered to the cap on Line~32.

We now focus on $M_2$. It changes what a withdrawal costs, and the test runs on
a path where it does so, yet the run itself shows nothing. A single run affords
only one kind of assertion, a copy of a value it produced, and here every such
value agrees with the reference. What $M_2$ breaks is the relation between the
amount withdrawn and the fee charged, and a relation cannot be stated at one
input. A natural thought is to enhance the existing test. Specifically, the
input values can be widened into a region that stays on the same path, and an
assertion can then be formed over that region.

Listing~\ref{lst:motivation-put} takes both steps. The deposit and the amount
withdrawn become separate parameters (Line~11), bounded rather than pinned
(Lines~12--14). The surrounding contract and the call sequence are unchanged
(Lines~1--9, 15--20). The last three lines are new. The fee is recovered from
the amount and what came back (Line~21). It is then required to lie between the
rates the bounded discount allows (Lines~22--23). The bounds keep every
admitted input on the path the counterexample took. A discount of at least one
basis point enters the conditional on Line~26 of Listing~\ref{lst:motivation}.
Deposits and amounts of this size keep the fee far under the cap on Line~32.
In the deposit and the amount the counterexample sits at a corner. Every
admitted amount must be at most every admitted deposit, or Line~24 would reject
some pair, and the counterexample withdraws its whole balance, so from there the
amount can only shrink and the deposit can only grow. The discount is not
constrained against either, and the region extends to both sides of the
counterexample's 39 basis points. The
bounds also give the assertion its content. A discount of one to a hundred
basis points leaves the contract charging a rate of 150 to 249, so the fee is
bounded by the amount times those two rates. $M_2$ charges on the deposit
instead of on the amount, and where the deposit is large enough against the
amount the fee leaves that range. The fee
it charges stays below every admitted amount, so the call completes and the
assertion is what reports the fault. $M_1$ was already reported by the concrete
test, and widening does not lose it. What widening adds is $M_2$.

These observations motivate \tool{VeriPUT}. A test is generated for a complete
path rather than for a condition, so the decisions the call takes are fixed and
an assertion has something to range over. Its arguments are not pinned to the
witness that reached the path but widened to a region of inputs that reach it
too. The assertion is synthesized over what the call changed, and the verifier
that produced the witness proves both the region and the assertion.

\section{Methodology}\label{sec:method}

\subsection{Overview}\label{sec:overview}

\begin{figure*}[t]
  \centering
  \includegraphics[width=\textwidth]{flowchart.png}
\caption{Overview of \tool{VeriPUT}, following one execution path through path
enumeration, certified input regions and verified assertion synthesis. All three
stages put their queries to the same verifier, drawn as the band above them. A
path the verifier cannot reach, or one whose input describes an execution the
contract cannot perform, emits \emph{no test} and is logged on its own; a path
whose region cannot be certified keeps its \emph{concrete replay test}; the rest
reach one or more \emph{parameterized unit tests}.}
  \label{fig:overview}
\end{figure*}

Generating a parameterized unit test for a contract calls for two things at
once: a set of inputs the test admits, and an assertion that holds for all of
them. Neither can be settled on its own. A set of inputs chosen first is usually
wide enough that the only assertions holding over it are ones that would hold
whatever the contract did. An assertion chosen first leaves the inputs on which
it holds still to be found, and those inputs need not form anything a test can
state as a range.

An execution path settles the two together. A path here is the sequence of
decisions one call takes, as recorded at the chosen granularity, from the entry
point to the point where it returns or reverts. An input is everything that
execution depends on: the arguments the call carries, the context it arrives in
such as the caller and the value sent, the state the contract is already in
such as a stored balance or owner, and what the contracts it calls out to
return. The inputs that follow a given path are those satisfying the conditions
along it, and all of them end at the same exit, so there is something definite
for an assertion to be about. \tool{VeriPUT} works in that order. As shown in
Figure~\ref{fig:overview}, it proceeds in three stages: \emph{path
enumeration}, which finds paths and one input for each; \emph{certified input
regions}, which widens that input into a range of inputs following the same
path; and \emph{verified assertion synthesis}, which proves assertions over the
range.

In the path enumeration stage, \tool{VeriPUT} instruments each entry point so
that the way a run went is appended to a record of its own at every decision
point the chosen granularity keeps, and plants at each exit an assertion denying
the path that reaches it. The record is written but never read back, so the
instrumented contract takes the same decisions as the original on every input.
Refuting one of those assertions hands back an input that reaches the path, and
from each such input a concrete replay test is emitted. Three kinds of path are
left without one and recorded apart: those the verifier proves out of reach,
those it leaves open, and those whose input describes an execution the contract
cannot perform.

The single input from the first stage says nothing about the other inputs that
reach the same path. In the certified input regions stage it is widened into a
range. No input takes two paths, so the other paths through the same entry point
propose cuts that keep the input already in hand, and the verifier is asked
whether the range that survives admits inputs of that path only. A refutation
supplies an input inside the range that left by a different path, which narrows
the range before the question is asked again. Where narrowing a coordinate would
leave it a single value, that coordinate is fixed at the value already in hand
and the others carry on, which is the common outcome. Where nothing survives,
the concrete replay test for the same path stands.

In the verified assertion synthesis stage, \tool{VeriPUT} supplies the oracle.
Candidates are enumerated from what the contract declares---its state variables,
the values it returns, and whether a call reverts---rather than predicted from a
model of what may go wrong, and each is put to the verifier with the certified
range as its assumption. Those it discharges are what the test carries. A candidate
refuted on part of the range is kept over the rest by dividing the range at the
input the verifier returned, and one that a condition separates is written with
that condition in front of it; a declined candidate is recorded as such. The
stage emits a parameterized unit test for each part. Where the verifier
discharges nothing over the range, the concrete replay test stands, asserting at
its single input the values the refutation already reported.

\paragraph{Why not write the test directly?}
Property inference suggests a shorter route: examine the contract, obtain a
property that holds of it, and write that property out as a test. The route does
not arrive at the same place. How much a property says is decided by the set of
runs it is asserted over, and that set is fixed by where the property is stated.

Take the withdrawal of Listing~\ref{lst:motivation}. A postcondition of it as a
whole has to hold on the runs that apply a discount and on the runs that do not,
so it can say no more than that the rate charged is at most the standard one. An
assertion of that kind is \emph{trivial}: it holds however the function behaves,
and a test carrying it reports nothing. Confine the question to one path and
each becomes worth checking. Along the path that applies a discount the rate is
strictly below the standard one; along the other it equals it. Both are
provable.

What changed is not the assertion but the set of runs it speaks about. Strength
comes from narrowing that set. Confining the question to one path leaves runs
that share a recorded decision sequence and an exit, and leaves every input in
exactly one class; the instrumented contract yields both the classes and their
exits. Beginning from a property reverses the problem, since what must
then be found is the set of inputs on which that property holds, and such a set
need not be anything an assumption can state. A parameterized unit test needs
both halves. Fixing the path first yields a set of inputs that can be written
down, and the assertion is proved over that set instead of sampled on it.

\subsubsection{Problem Definition}\label{sec:task-def}

The input to \tool{VeriPUT} is the source of a single contract and nothing
else: no test written by a developer, no property supplied by hand, and no
record of transactions the contract has served. For each path of each
\emph{unit}, that is, each entry point the outside world may call, it produces
one of three outcomes: one or more parameterized unit tests, each admitting a
part of the certified range on which its assertions were proved; a concrete
replay test, where no range could be certified or where nothing the test writes
was left free to vary; or no test at all, the path being listed with the
reason. A constructor runs only when the contract is created and is not an
entry point of the deployed contract, so its execution covers no path of a
unit. Every test carries at least one assertion, proved to hold on each input
the test admits and not only on the values a framework happens to draw.

Results are relative to bounds we declare: how far loops are unrolled and
callbacks followed, how long an argument of dynamic length may be, how far
internal calls are expanded, and how many values a coordinate's range may have
removed from it one at a time. We assume a decoder that rejects arguments not
fitting the declared types instead of normalising them. Gas consumption is not
modelled, though the branches introduced by a gas failure are. Each run of
\tool{VeriPUT} also fixes a \emph{fixture}, a deployment a test can reproduce:
the code, the values construction leaves behind, and the block context. The
state a call starts from is the state that fixture leaves, except at the
storage locations a test can write, which range as the other inputs do. The
fixture is settled once, before enumeration. \tool{VeriPUT} asks the verifier
for one deployment that completes, fills the values that witness leaves
unconstrained the same way every time, and runs that deployment once to see it
through; every question below is put from what it leaves behind, and the emitted
test repeats the same deployment. Where no deployment
can be witnessed within these bounds, the contract yields no tests and is
reported. A path is reachable relative to that fixture and to these bounds, and
not from every deployment or every history of calls. The claims below concern
the queries put to the verifier and take the verifier to be free of defects;
where it exhausts its resources it says so, which we read as an answer withheld
and not as a proof.

Four things are asked of the verifier itself. (i)~It decides several assertions
under one assumption in a single run, which is what keeps a ladder of candidate
bounds to one round. (ii)~It returns, with a refutation, the values that produce
it. (iii)~It reports a question it could not settle as such, and not as one that
holds. (iv)~It reports the value a variable of the instrumented unit carries,
which is how a refutation names the path it took. A verifier missing one of
these leaves the construction sound and weakens it. Without the first,
\tool{VeriPUT} is run with the ladder switched off.

\subsection{Path Enumeration}\label{sec:enum}

We now describe how \tool{VeriPUT} obtains the paths a parameterized unit test
will later be built around, together with one input for each. Recall from
Section~\ref{sec:task-def} that a \emph{unit} $u$ is one entry point of the
contract, and that a test for a path of $u$ has to state which inputs it admits
and assert something about what the call left behind. Both halves have to be
readable from the outside. A test written for a framework cannot assert that a
particular branch was taken, since a branch is a moment inside the call and
leaves no trace once the call has finished; it can only assert over what the
call returned and what it changed. The coverage target therefore has to fix
those, and a whole path does.

\paragraph{Why must the coverage object be a whole path?}
\begin{itemize}
\item A target inside a run leaves the ending open. Two inputs may take the same
branch and then part company, one returning and one reverting, so a test bound
to that branch has nothing definite to check when the call is over.
\item The targets have to cut the inputs into classes a test can name. The
inputs reaching a branch are those reaching it along any route, and one input
reaches many branches. Whole paths leave each input in exactly one class, which
is what Section~\ref{sec:region} relies on when it bounds one class by the
others.
\end{itemize}

\subsubsection{Decision Points}
A \emph{decision point} of a unit is a place where its control flow may continue
in more than one way. They come from four sources. The first is the
guard of a conditional or a loop. The second is the right-hand operand of a
short-circuit operator: in \texttt{a \&\& b} the operand \texttt{b} is evaluated
only when \texttt{a} holds, so a run in which \texttt{a} fails and a run in
which \texttt{a} holds while \texttt{b} fails are different executions, one of
which evaluated \texttt{b} and one of which did not. The third is the outcome of
a call to another contract, which may succeed or fail, and, where the called
contract calls back, the entry point it calls back into. The fourth is the
failure of argument decoding, which ends the call before its body starts.

Notably, the checks a compiler inserts for arithmetic and for array bounds are
not counted among these. The line falls between what the source states and what
a particular compiler adds. Drawing it there keeps one source to one set of
decision points across compiler versions, and across blocks in which those
checks are switched off. Inputs that would trip an inserted check are removed
from a range at the point where the range is certified, as
Section~\ref{sec:region} describes. A path never tells
an arithmetic failure apart from the operation succeeding, and the number of
paths through a unit is a property of the code its author wrote.

Which decision points are recorded is a parameter. Expanding an internal call
inside its caller multiplies the decisions of the one by the decisions of the
other, so \tool{VeriPUT} records those of a chosen subset of the internal call
sites of $u$. We write $S(u)$ for those call sites, $G \subseteq S(u)$ for the
subset that is expanded, which we call the \emph{granularity}, $\Delta_G(u)$ for
the decision points recorded under it, and $\Pi_G(u)$ for the resulting paths.
Section~\ref{sec:budget} returns to how $G$ is chosen.

\subsubsection{Instrumentation}
We then add two things to $u$. The first is a record: at each point of
$\Delta_G(u)$, the way the run went is appended to a variable declared for this
purpose. The record is written and never read back, and nothing already in the
contract reads it either, so the instrumented unit takes the same decisions as
the original on every input. The second is one assertion per \emph{exit} per
enumerated path. A unit has two kinds of exit, one at which the call returns and
one at which it reverts, and the assertion planted at an exit claims that the
path in question was not the one just taken. An assertion of this form is false
exactly on the runs that do take the path, so asking a verifier to prove it is a
way of asking whether any input reaches it, and a refutation carries such an
input back.

Where the assertion is placed depends on which kind of exit it belongs to. At an
exit that returns, it goes at the end of the body. At an exit that reverts, it
goes ahead of the operation that fails, because a revert undoes everything the
call wrote and nothing after that operation runs. Placing it after the failure
would leave it unreachable on every input, and a verifier would report that it
holds, which reads exactly like a proof that the path cannot be taken.

Listing~\ref{lst:instrumented} shows the withdrawal of
Listing~\ref{lst:motivation} instrumented for two of its paths, one leaving by
each kind of exit. The record is kept as an integer: Line~4 seeds it for the
call, and Lines~6, 11 and~15 double it and add the way the decision went, so
that a path is named by the number its decisions build up. Seeding with one
rather than zero keeps a run that has taken no decision apart from one that has
taken a single decision the low way. Line~7 then targets the path on which the
amount is rejected, and Line~18 the path on which a discount applies and the fee
stays under the cap. A refutation of the second returns the discounted
withdrawal at the corner of the region Listing~\ref{lst:motivation-put} draws.

The two assertions sit at different places, and the rule above is why. The
first path leaves by reverting, so its assertion goes ahead of the operation
that fails, an assertion being checked where it stands and before the undoing.
The second returns, so its assertion goes at the end of the body. Placed
earlier, it would be refuted by a run that later reverts in a transfer, and the
input handed back would reach an exit the path does not name.

The listing appears in source form for readability, whereas the instrumentation
is applied to the control-flow graph, where a conditional expression and a
short-circuit operator have already become ordinary branches; and the record is
a quantity of the model rather than a storage variable of the contract, so what
a revert does to the contract's own state does not reach it. The record itself
is the sequence of ways the run went; the integer of the listing is one
encoding of it, chosen so the example reads, and it carries no bound of its
own. We also reduce the example in four ways. (i)~Each condition is lifted into
a named local, so the record and the branch it belongs to read on separate
lines. (ii)~The two operands of the conjunction at Line~5 are shown as one
decision, where the encoding treats them as two. (iii)~The two transfers that
end the function are left folded, each carrying decisions of its own. (iv)~The
decision the decoder takes before the body starts is left out.

\begin{lstlisting}[language=Solidity, numbers=left, numberstyle=\scriptsize,
xleftmargin=1.5em, basicstyle=\ttfamily\scriptsize, breaklines=true,
columns=fullflexible, keepspaces=true, float=tbp,
caption={The withdrawal of Listing~\ref{lst:motivation} instrumented for two of
its paths, one leaving by each kind of exit.},
label={lst:instrumented}]
uint256 tr;                       // decision record

function withdraw(uint256 amount) public returns (uint256 net) {
    tr = 1;                       // fresh for this call
    bool ok = amount > 0 && amount <= deposits[msg.sender];
    tr = 2 * tr + (ok ? 1 : 0);
    assert(tr != 2);              // target: the amount is rejected
    require(ok, "bad amount");
    uint256 rate = feeBps;
    bool disc = discountBps[msg.sender] > 0;
    tr = 2 * tr + (disc ? 1 : 0);
    if (disc) { rate = feeBps - discountBps[msg.sender]; }
    uint256 fee = amount * rate / 10000;
    bool cap = fee > maxFee;
    tr = 2 * tr + (cap ? 1 : 0);
    if (cap) { fee = maxFee; }
    // ... two transfers and the event, each with decisions of its own
    assert(tr != 14);             // target: discounted, under the cap
}
\end{lstlisting}

\subsubsection{Verdicts}
What \tool{VeriPUT} hands the verifier is the instrumented unit together with
the assertions planted in it. What the verifier hands back, for each assertion,
is an answer and, where that answer is a refutation, an input producing it. The
answer comes in one of three kinds, and \tool{VeriPUT} keeps the three apart. It
may refute the assertion, in which case the path is reachable and an input
reaching it is in hand. It may prove the assertion, in which case no input
reaches the path. It may run out of resources without deciding, in which case
nothing is known about the path, and it is reported as such and never counted
with either of the others.

Two conditions have to hold for the second kind of answer to mean what it says.
The first concerns the reasoning behind it. A bounded model checker unrolls
loops and follows calls to a fixed depth, and an assertion it cannot refute
within those limits has been shown to hold only of the executions the limits
admit. Read as a statement about the contract, such an answer belongs with the
third kind and not the second. Only where the verifier settles the question
without leaning on a bound, as unbounded reasoning about the unit can, does the
answer say that no input reaches the path. Therefore \tool{VeriPUT} counts a
path as unreachable only when the answer is of the latter sort.

The second condition concerns what was reasoned about. A path guarded by state
that only an earlier call can establish is reached by writing that state, as
Section~\ref{sec:task-def} declares, and not by modelling the call that would
have written it. What a refutation witnesses is a path of the fixture it was
asked about. Moreover, where the
reachability of a path rests on a construct the model describes more loosely
than the contract does, that path is set aside before any test is built from it,
and it is listed together with the construct responsible. Unlike a path left
undecided, where the answer is missing, here the answer would be about an
execution the contract cannot perform.

\subsubsection{Budget and Granularity}\label{sec:budget}
Neither response below removes a path from the contract. One folds paths
together and one stops enumerating them, and the record is written on every run
either way, so nothing the later stages do rests on the paths having been
enumerated in full. The responses decide how much is produced, never whether it
is correct.

A unit whose callees are themselves layered can carry more paths than any budget
allows. \tool{VeriPUT} responds by withdrawing call sites from $G$. A withdrawn
call still executes and still determines what the run does; what it stops doing
is contributing its own decisions to the identity of a path. Paths that differed
only inside it become one path, which still runs to an exit and still holds a
class of inputs of its own. Therefore every input remains in exactly one class,
and the loss is confined to how finely the classes are told apart.

The oracle pays that loss, and the region does not. Section~\ref{sec:oracle}
loses the statements it reads off a path, for the decisions no longer recorded.
It keeps the ones comparing an observable before and after the call, at the
price of the further tests a division costs. That section returns to both.

Given a budget $N$ on the number of targets, the subset $G$ is chosen in three
steps.

\begin{enumerate}
\item \textbf{Measurement, not estimation.} For each candidate site $s$,
\tool{VeriPUT} expands every site but $s$ and counts the resulting paths, which
gives a number $r(s)$: how many paths remain when $s$ alone is withdrawn. A
small $r(s)$ marks a site that was holding up many paths. The count stops at the
budget, so every site that cannot on its own bring the unit under it receives
the same $r$ and the step does not separate them.
\item \textbf{Greedy withdrawal.} The candidates are ordered once by increasing
$r$, and sites are taken from the front of that order and withdrawn until the
count of paths under the remaining $G$ falls to $N$ or below.
\item \textbf{Distance.} Sites with equal $r$ are separated in favour of the one
furthest from the entry. The reasoning is a heuristic: decisions taken close to
the entry are the ones most directly determined by the arguments of the call,
and keeping them inside the identity of a path keeps the classes aligned with
what a test can vary. On a unit far over the budget the first step separates
few sites, so most of the order comes from here.
\end{enumerate}

The order is settled once and never recomputed, since withdrawing one site can
make the measurement of another stale and recomputing would cost one count per
remaining site per round. We do not claim the result to be optimal, and none of
the guarantees above depends on which sites are chosen. The choice decides how
much strength is kept, never whether the output is sound.

One class of call site is never a candidate. Where the callee of a site is
itself an entry point, withdrawing the site leaves the call arriving at a body
that carries the guard the compiler places on external entry, a guard an
internal call never passes on chain. The model would then admit a failure the
contract cannot produce, and a test built from it would fail on the contract as
shipped.

Where withdrawal is not enough, enumeration stops at $N$ targets and the paths
beyond it are counted and reported. Doing so removes assertions and leaves the
record untouched, so an input taking one of the abandoned paths still carries
that path in the record when the call ends. Section~\ref{sec:region} returns to
this point, which is what allows the later stages to remain correct without
every path having been enumerated.

Algorithm~\ref{alg:instrument} collects the stage. It reads the internal call
sites of $u$, sets aside those whose callee is an entry point, measures and
withdraws as above, counts and enumerates the paths of $\Pi_G(u)$ up to $N$
under the declared unrolling and granularity, pruning prefixes the control flow
cannot follow, appends the record at each point of $\Delta_G(u)$,
and plants one assertion per path at the exit that path reaches.

\begin{algorithm}[!htpb]
\caption{Instrumenting one unit for path enumeration}
\label{alg:instrument}
\begin{algorithmic}[1]
\Require unit $u$, target budget $N$
\Ensure instrumented unit, granularity $G$, abandoned path count $k$

\State $G \gets \textsc{InternalCallSites}(u)$
\Comment{those reached within the declared call depth}
\State $E \gets \{s \in G : \textsc{Callee}(s)$ is an entry point$\}$
\Comment{never candidates}
\If{$\textsc{CountSyntacticPaths}(u, G) > N$}
    \ForAll{$s \in G \setminus E$}
        \State $r(s) \gets \textsc{CountSyntacticPaths}(u, G \setminus \{s\})$
        \Comment{one count per site}
    \EndFor
    \State $Q \gets G \setminus E$ ordered by increasing $r$,
           ties by distance from the entry
    \While{$\textsc{CountSyntacticPaths}(u, G) > N$ \textbf{and} $Q \neq \emptyset$}
        \State remove the head of $Q$ from $G$
        \Comment{the order is never recomputed}
    \EndWhile
\EndIf

\State $\Pi \gets \textsc{EnumeratePaths}(u, G, N)$
\Comment{structurally impossible control-flow prefixes are pruned}
\State $k \gets \textsc{CountSyntacticPaths}(u, G) - |\Pi|$

\ForAll{$p \in \Delta_G(u)$}
    \State $\textsc{AppendToRecord}(p)$
    \Comment{written, never read}
\EndFor
\ForAll{$\pi \in \Pi$}
    \If{the exit of $\pi$ reverts}
        \State $\textsc{PlantAssert}(\mathit{record} \neq \pi)$
               before the operation that fails
    \Else
        \State $\textsc{PlantAssert}(\mathit{record} \neq \pi)$ at that exit
    \EndIf
\EndFor
\State \Return the instrumented unit, $G$, $k$
\end{algorithmic}
\end{algorithm}

\subsubsection{Concrete Replay Tests}
Each refutation yields an input, and from it \tool{VeriPUT} emits a test that
sets every argument, the caller, the value sent, and the contract state to the
values the input names, and then calls $u$. The test is bound to the path the
input takes, which is what allows the assertions of Section~\ref{sec:oracle} to
be attached to it without further queries. The refutation also reports the value
each observable holds when the call ends, so such a test can state those values
without a query of its own. Moreover, it gives every path with a
witness a test of its own before the widening of Section~\ref{sec:region} is
attempted.

Some of the values an input carries stand for nothing in particular. A verifier
with no reason to constrain an address may report it as zero and use it only to
tell one account from another. Such a value marks a place where the path imposes
no condition, and the test is free to use any value the framework allows. Where
a value cannot be arranged at all, no test is emitted, and the path is reported
with the others that received none.

\subsection{Certified Input Regions}\label{sec:region}

Section~\ref{sec:enum} leaves one input for each path it witnesses. That input
shows the path can be taken and says nothing about the others that take it, and
those others are where a change in behaviour would show. We now widen it. We
write $D_\pi$ for the \emph{path domain} of $\pi$: the modelled inputs, under
the declared bounds, whose run records $\pi$ at the granularity of
Section~\ref{sec:budget}. This stage receives a path $\pi$, the witness $x_\pi$
that reaches it, and the witnesses of the other paths through the same unit,
and it returns a region of inputs proven to follow $\pi$, obtained by taking
what those other paths cover away from what $\pi$ might cover. Where no region
survives, the concrete replay test of the previous stage stands.

\subsubsection{Coordinates}\label{sec:coords}
A \emph{coordinate} is a dimension of the input that the emitted test can set
to a value of its choosing, directly or through its preamble. The arguments of
the call are set by passing them. The caller and the value sent are set through
the interfaces a test framework provides for the purpose. The state the
contract is in is set by writing it. The behaviour of a contract called out to
is set by deploying a stub that answers as the input names, which presumes that
the callee's interface is known and that a stub can stand in for it; where the
address is one the test cannot control, or no code is deployed there when the
call is made, that behaviour is not settable. Which behaviours are settable is a
parameter of the method: a behaviour counts when the verifier models it and the
emitted test can reproduce it, and the two must agree.

Which dimensions there are is settled by the types that may cross a contract's
boundary, and every one of them comes apart into integers. An enumeration
becomes a small unsigned integer, and a user-defined value type the type beneath
it. A handle to a contract or an interface becomes a $160$-bit address, and a
boolean an eight-bit scalar of which two values are used. A struct or an array
of fixed length gives one coordinate per component. An argument of dynamic
length gives one coordinate for its length and one per element it may hold, up
to the bound declared in Section~\ref{sec:task-def}. The decomposition terminates without a
depth bound of its own, because a type that could recur without end, a mapping
or a pointer into storage among them, cannot cross the boundary at all.

The state a call starts from decomposes by storage location instead. A path
fixes which locations it reads and writes wherever its decisions are recorded,
and such a location becomes a coordinate where its key is an expression over the
input and the test can establish it from the values it renders; the deposit of
the caller in Listing~\ref{lst:motivation} is one. Locations reached inside a
withdrawn call site, locations whose keys can coincide, and locations no test
can write receive none.

Every coordinate ranges over a bounded set of integers, so there is no kind of
argument the construction below is unable to describe. What a given
implementation admits is narrower: the kinds for which the verifier has a model
and the emission of Section~\ref{sec:emit} has a way of writing a value, and a
quantity of any other kind is reported rather than approximated. How much of a
coordinate a region ends up admitting is a quantity we report rather than a
capability we claim.

What a region can say about a coordinate depends on the type beneath it, and
three cases arise. (i)~On a quantity that counts or measures, a range carries
information at every width. (ii)~On an address it does not. A path turns on
which account the caller is, and not on where that account falls among the
addresses. A range of addresses narrower than the type and wider than one value
is therefore sound and still leaves a test varying a quantity the unit never
orders. The brackets and cuts of Section~\ref{sec:candidate} produced that
range, so \tool{VeriPUT} emits it; Section~\ref{sec:emit} declines to count it.
A boolean is the same case with two values. (iii)~On an argument of dynamic
length, the length is one coordinate and each element it may hold is another. An
element the length excludes is never read, so the region leaves those elements
whole and Section~\ref{sec:emit} renders the length first.

Not every quantity a verifier fixes is a coordinate. A counterexample names the
quantities it chose, and each is classified by kind: an argument, part of the
context, a slot of the contract's state, or a value returned to the unit by a
contract it called. A quantity is a coordinate when the emission of
Section~\ref{sec:emit} has a way of establishing that kind, and the
classification is closed, a kind it does not recognise counting as not settable.
Where a quantity that decides the path is not a coordinate, $\pi$ receives no
region and no test, and is reported with the quantity responsible. Taking a
settable quantity for an unsettable one costs yield; the opposite mistake is
caught at run time, as Section~\ref{sec:emit} describes. A quantity
left out of the classification altogether is not constrained by the region
either, so the certification of Section~\ref{sec:cert} quantifies over it: where
varying it can lead away from $\pi$, that query is refuted or left undecided and
no region beyond the witness is emitted.

\paragraph{Why not widen the witness one step at a time?}
One alternative widens the witness a little, asks whether the widened range
still follows $\pi$, and repeats. Taken one step at a time that is a search with
no stopping rule, in which a step too small spends a query and a step too large
is refuted, while neither outcome says how far the boundary lies. The
construction below starts from the witness as well, but fixes the widths in
advance and puts them in one batch, so the answers together bracket the boundary
instead of feeling for it. The brackets of the other paths through $u$ are asked in the same batch, so
cutting against them adds no further round of verification. They offer candidate
cuts and not exact boundaries: projections of disjoint domains can overlap, and
where they cover a coordinate entirely nothing is left and the concrete replay
test stands.

\subsubsection{Candidate Regions}\label{sec:candidate}
A \emph{certified input region} for $\pi$ admits a set of values on each
coordinate and imposes nothing across coordinates. Formally,
\begin{equation}
R_\pi \;=\; \prod_{c \in \mathrm{coords}(u)} V_c ,
\qquad
V_c \;=\; [L_c, U_c] \setminus H_c ,
\label{eq:region}
\end{equation}
where $[L_c,U_c]$ is an interval of the integers $c$ ranges over and $H_c$ is a
finite set of \emph{punctures}, values that interval contains and the region
excludes. The steps below show where punctures come from and why intervals alone
are not enough.

Imposing nothing across coordinates is a decision about what is inferred. The
construction below reads each coordinate on its own, so it derives no relation
between two of them for a test to carry. A path domain need not be a product,
however. On the path of Listing~\ref{lst:motivation} where a discount applies
and the fee stays under the cap, the two are tied by
$\lfloor \mathit{amount} \times (\mathit{feeBps} - \mathit{disc}) / 10^{4}
\rfloor \leq \mathit{maxFee}$, which no pair of independent ranges describes.
The same condition gives the cost. At the smallest discount the
path admits, the fee reaches the cap at about four ether. Where the discount
leaves a single basis point of rate in force, the same cap is reached only at
two hundred and forty-nine times that amount, and where it leaves none the cap
is never reached at all. A region giving the amount one range and the discount
another has to take the smallest of these, so a withdrawal of ten ether by an
account holding a large discount follows this path and receives no test from it.
We therefore make no claim that a certified region is all of a path domain, and
what one leaves out is counted rather than passed over.

A candidate region is built in three steps, of which only the first and third
put questions to the verifier.

\begin{enumerate}
\item \textbf{Bracketing.} For each coordinate we ask whether it stays below a
candidate bound on every run that follows $\pi$, and again whether it stays
above one. The candidates form a ladder anchored at the witness:
$x_{\pi,c}$ itself, then $x_{\pi,c} \pm 1$, $\pm 2$, $\pm 4$ and so on to the
extremes of the coordinate's own type. The spacing doubles because on a
coordinate of $256$ bits any evenly spaced ladder puts all of its rungs where
nothing is. The anchor is the witness because the witness already answers one
side of each question. A bound below $x_{\pi,c}$ is refuted by $x_\pi$ itself,
and a ladder anchored at zero spends most of its rungs on answers of that kind.
One rung is still placed below $x_{\pi,c}$, where a refutation is what we
expect. A discharge there means that no run following $\pi$ satisfies the
assumption the batch was asked under, and that every answer in it holds for want
of a run to violate it. Without that rung, a coordinate fixed at one value and
an assumption nothing satisfies are reported alike. The candidates do not depend on one another, so a verifier that
decides several assertions under one assumption settles a ladder at once, and
the ladder is not walked rung by rung. What comes back is an
over-approximation: an interval known to contain everything $D_\pi$ reaches on
that coordinate. A candidate the verifier leaves undecided is dropped, which
widens the bracket.

\item \textbf{Subtraction.} No input takes two paths, so a value another path
through $u$ reaches on every coordinate at once cannot be $\pi$'s, and the
brackets of those paths are where we look for cuts. A bracket lying wholly to
one side of $x_{\pi,c}$ cuts the interval there; one that contains $x_{\pi,c}$
cuts nothing on that coordinate. Because projections of disjoint domains can
overlap, a cut may take away values $\pi$ does reach, which costs yield alone,
and may leave values it does not, which certification removes. Subtraction
proposes the candidate; certification decides it. Where enumeration was cut
short, fewer sibling brackets are available and the candidate comes out wider.

\item \textbf{Punching.} A sibling that pins a coordinate to one value needs
care, since an interval cannot say \emph{everything except} $v$. Whether a
sibling $\sigma$ is so pinned on a coordinate $c$ is answered by running the
first step on the single candidate $v$, where the two probes ask whether $c$ is
at once at most and at least $v$; we call this a \emph{collapse test}. The
candidate costs nothing to obtain, being the value $\sigma$'s own witness gives
to $c$, and it goes into the same batch. A discharged collapse test says that
$\sigma$ occupies exactly $\{v\}$ on $c$, and where $v$ differs from
$x_{\pi,c}$ it joins $H_c$ as a puncture instead of cutting $[L_c, U_c]$ in
two.
\end{enumerate}

The guard on Line~18 of Listing~\ref{lst:motivation} is the shape this handles.
Take the path of \texttt{setDiscount} on which the guard fails. Its sibling, the
path on which the guard holds, is reached only where the caller is the owner,
one value somewhere inside the address space. Subtraction confined to intervals
would have to keep one side of that value and discard the other, and which side
it kept would follow from the value the solver happened to return: two answers,
both sound, one holding almost the whole space and the other almost none of it.
With a puncture the region is the address space less that one value, wherever it
sits.

A puncture answers that shape and no other. Where two coordinates are tied
together, as the amount and the discount are by the cap, no product of value
sets describes the domain however many values are removed from it, and the
shortfall is reported as yield rather than repaired.

Algorithm~\ref{alg:region} collects the three steps together with the
refinement of the next section. It brackets each coordinate of $\pi$, subtracts
the brackets of the other paths through $u$ on the side away from $x_\pi$,
asks a collapse test at each sibling's own witness value and punches out the
ones it pins, puts the result to the verifier,
and, on a refutation, narrows the region at the input the verifier returns
before asking again. That input is punched out where it differs from the witness
on one coordinate and the budget on holes allows, cut against otherwise, and
where a cut would leave the coordinate it points at a single value, that
coordinate is pinned to the value $x_\pi$ gives it and the others carry on. Rounds that run out return the region as it then stands, the witness
alone only where every coordinate is pinned; a query the verifier declines
returns the witness alone, there being nothing to cut at.

\begin{algorithm}[!htpb]
\caption{Building a certified input region for one path}
\label{alg:region}
\begin{algorithmic}[1]
\Require unit $u$, path $\pi$, its witness $x_\pi$, the enumerated paths $\Sigma$
         through $u$ that carry witnesses, refinement rounds $m$, hole budget $h$
\Ensure a certified region $R_\pi \subseteq D_\pi$

\ForAll{$c \in \mathrm{coords}(u)$}
    \State $V_c \gets \textsc{Bracket}(u, \pi, c)$
    \Comment{one ladder, over-approximating}
    \ForAll{$\sigma \in \Sigma$}
        \State $v \gets$ the value $\sigma$'s witness gives $c$
        \If{$v \neq x_{\pi,c}$ \textbf{and} $\textsc{Collapse}(u, \sigma, c, v)$}
            \State $V_c \gets V_c \setminus \{v\}$
            \Comment{$\sigma$ occupies $\{v\}$ here; $v$ becomes a puncture}
        \Else
            \State $V_c \gets \textsc{CutTowards}(V_c,
                   \textsc{Bracket}(u, \sigma, c), x_{\pi,c})$
            \Comment{no cut where the bracket holds $x_{\pi,c}$}
        \EndIf
    \EndFor
\EndFor
\State $R \gets \prod_{c} V_c$

\While{$m > 0$}
    \If{\textbf{not} $\textsc{KnownNonEmpty}(R)$}
        \Comment{an empty region is never reported as certified}
        \State \Return $\{x_\pi\}$
    \EndIf
    \State $(\mathit{verdict}, y) \gets \textsc{Certify}(u, R, \pi)$
    \If{$\mathit{verdict}$ is discharged}
        \State \Return $R$
    \ElsIf{$\mathit{verdict}$ is declined}
        \State \Return $\{x_\pi\}$
        \Comment{no counterexample to cut at}
    \EndIf
    \If{$y$ differs from $x_\pi$ on one coordinate $c$ \textbf{and} $|H_c| < h$}
        \State $R \gets \textsc{Punch}(R, c, y_c)$
        \Comment{removes $y$ alone; the next round decides the rest}
    \ElsIf{some coordinate admits a cut at $y$ leaving it more than one value}
        \State $R \gets \textsc{Cut}(R, y, x_\pi)$
        \Comment{the query failed on $y$, not on $x_\pi$}
    \Else
        \State $R \gets \textsc{Pin}(R, y, x_\pi)$
        \Comment{fix that coordinate at the witness; the others carry on}
    \EndIf
    \State $m \gets m - 1$
\EndWhile
\State \Return $R$
\Comment{pinned where cutting could not narrow; $\{x_\pi\}$ where all are}
\end{algorithmic}
\end{algorithm}

\subsubsection{Certification}\label{sec:cert}
We now put one question to the verifier. In words: assuming the inputs lie in
this region, does the call take $\pi$ and leave by the exit $\pi$ reaches?
Formally, $R_\pi$ is supplied as the assumption, and an assertion stating that
the run reaching this exit took $\pi$ is planted at every exit of $u$.

At every exit, and not only at the exit $\pi$ reaches, for two reasons.
(i)~An input inside the region that leaves along another path leaves by another
exit, and were the assertion planted only where $\pi$ ends it would never meet
one, leaving the query to hold for want of anything able to refute it. (ii)~An
input may carry the record of $\pi$ and still leave elsewhere, where a decision
the granularity does not record ends the call early; at that exit no run can
satisfy the assertion, so such an input refutes it. Without the second, a region
could be certified for a path whose exit it does not fix.

Three answers are possible, and \tool{VeriPUT} keeps them apart. The verifier may
discharge the query, and the region is certified. It may refute it and return an
input inside the region on which one of the planted assertions fails, either
because that input left along another path or because it tripped an inserted
check; that input, together with $x_\pi$, which is known to lie in $D_\pi$, gives
a definite cut, and the region is narrowed and the question asked again. Or the
query may be declined.

A refutation may differ from $x_\pi$ on several coordinates at once. On any
coordinate where they do differ, keeping
the side of $y_c$ on which $x_{\pi,c}$ lies removes $y$ and keeps $x_\pi$, so
every such coordinate offers a sound cut, and \tool{VeriPUT} takes the one that
leaves the largest fraction of its coordinate's values. A region is a product,
so a cut costs the fraction it removes and not the count. A cut that takes
fifty values from a hundred costs more than one that takes nine hundred from a
thousand.
Where $y$ and $x_\pi$ agree on every coordinate, they differ only on a quantity
no test can set, and the path goes to the gate of Section~\ref{sec:coords}
instead of to another round.

Where $y$ differs from $x_\pi$ on one coordinate alone, removing that single
value excludes $y$ and keeps every other value the region held, which no cut
does. \tool{VeriPUT} punches it out instead, up to the budget on holes declared
in Section~\ref{sec:task-def}. Such a puncture differs from those of
Section~\ref{sec:candidate}, which a collapse test established. This one rests
on a single refuted input. It removes the input in hand and says nothing about
the values beside it, and the round that follows decides whether more remain.

An assumption no input satisfies makes every assertion under it hold for want of
a run to violate it, so such an outcome is indistinguishable, by its verdict
alone, from a region the verifier has checked. Two ways of building one arise
from our own construction: punching can empty an interval, and a bound written
for a coordinate may not fit that coordinate's type. Both are checked before the
question is put, and where satisfiability cannot be settled that way it is asked
on its own; that check answers yes only where non-emptiness has been
established, so an undecided answer and an empty region leave by the same exit.
A declined query is reported as declined and never as a certification.

The same hazard reaches every assumption the construction writes, and one rule
keeps it out of all of them. Every assumption is built from a witness that
satisfies it. The bracketing step asks about runs following $\pi$ under values
$x_\pi$ carries, and a collapse test asks about runs following $\sigma$ at the
value $\sigma$'s own witness gives. A rung whose expected answer is a refutation
coming back discharged is what shows the rule broken.

The same question carries the checks a compiler inserts. Section~\ref{sec:enum}
kept those out of the identity of a path; here they are turned on, so that an
input in the region that would overflow or index out of bounds refutes an
assertion of its own and is cut away with the machinery already in place. Each
is represented as an assertion, not as a recorded decision, so the number
of paths is unchanged. A discharged region therefore contains only inputs that
take the recorded decisions of the contract as written and satisfy every
inserted check.

Refinement need not end at a single point, and usually should not. Where a
refutation cannot be answered by cutting the coordinate it points at, or where
cutting would leave that coordinate one value, \tool{VeriPUT} pins the
coordinate to the value $x_\pi$ gives it and carries on with the others. A region
with some coordinates pinned and others left to range says less than one in which
all of them range, more than a test at a single point, and it is the common
outcome.

Where the rounds run out on every coordinate, $R_\pi = \{x_\pi\}$, and a region
of one point leaves nothing for a test to range over. That the witness
follows $\pi$ needs no argument, since the path domains are pairwise disjoint
and $x_\pi$ lies in $D_\pi$. The inserted checks are another matter, because the
enumeration of Section~\ref{sec:enum} left them out and a witness may trip one.
The same question is therefore put on the single point, where every quantity is
fixed and the answer is a matter of evaluation. A witness that fails it receives
no test and is reported with the path.

\subsubsection{From a Region to a Test}\label{sec:emit}
A region reaches a test as an assignment and not as a filter. In a testing
framework a precondition that does not hold is a failing test, not a draw
skipped, and the primitive that does skip a draw both throws the work away and
is capped in how often it may be used. We write $A_t$ for the region a test $t$
renders, which is $R_\pi$ itself unless the punctures below force it to be
drawn in pieces, and which is in either case a product of coordinate sets. Each
coordinate is therefore given a value inside its set by construction, and the
punctures are the one thing left to skipping, in a proportion bounded by their
number against the width of the interval. Where they are dense enough that
skipping would exhaust that allowance, the interval is drawn in puncture-free
pieces, each emitted as a test of its own and each carrying the assertions
proved over the whole of $R_\pi$, since all of them lie inside it. What the budget on tests for
one path leaves undrawn is reported.

Some coordinates need no rendering at all. What matters is whether the values
the test environment would supply for $c$ on its own already lie inside $V_c$.
It is not whether $V_c$ is narrower than the type of $c$: that reading holds
for an argument the framework draws over the whole type, and fails for the
state a call starts from, where the environment may supply a single value left
by deployment, one that can sit outside $V_c$ altogether.

What the environment supplies is not knowable in general, so \tool{VeriPUT}
makes it knowable. Every test opens with the same preamble, which puts the
environment into a state the omission rule can be stated against, and every test
checks, before it calls the unit, that each coordinate holds a value inside
$V_c$, the omitted ones included.

That check is not a filter under another name. A filter discards draws that lie
outside the region legitimately, which is ordinary and must not fail a test. The
check looks at coordinates the framework does not draw at all, coordinates the
preamble was supposed to have established. A value outside $V_c$ there means the
preamble did not do what the omission rule assumed of it, which is a defect on
our side and not an input to be discarded. It therefore fails loudly, and it is
what makes the omission rule falsifiable rather than asserted.

What the rendering leaves decides which of the two tests a path receives. A test
is parameterized when at least one coordinate it renders is left more than one
value to take and is a coordinate whose order the text of the unit uses. The
second condition is what keeps a range that carries nothing from counting: on a
quantity a unit only tests for equality, such as an account, a range wider than
one value and narrower than the type admits inputs the unit cannot tell apart
from one another. The range is emitted and reported; what it does not do is
decide which of the two tests the path receives. A region wider than a point does not settle this on its own,
since the coordinates the omission rule leaves out are not rendered and a region
can be wide only on those. Where no rendered coordinate takes more than one
value, what the path receives is the concrete replay test of
Section~\ref{sec:enum}, whatever the width of $R_\pi$.

\subsubsection{Soundness of Emission}\label{sec:sound-emit}
We now show that a test emitted for a path speaks only about that path: every
input it admits follows the path it names, and every assertion it carries holds
on that input. Four things are needed, and each is established elsewhere rather
than assumed here. (i) The record has to be written at every decision point,
which is what the instrumentation of Section~\ref{sec:enum} does. (ii) Every
coordinate the region constrains has to be established by the test itself,
which the preamble of Section~\ref{sec:emit} does and its check confirms while
the test runs. (iii) The emitted test has to deploy the fixture the queries
were put under and to answer its external calls as they assumed, which the same
preamble does. And (iv) a declined query must not be read as a discharged one,
which is what the gate of Section~\ref{sec:cert} is for.

\textbf{Proposition (Soundness of Emission).}
Let $t$ be a test emitted for a path $\pi$ of a unit $u$ from a certified region
$R_\pi$, and let (i) to (iv) above hold. Then, under the bounds declared in
Section~\ref{sec:task-def}, every input $t$ admits lies in
$A_t \subseteq R_\pi \subseteq D_\pi$, and every assertion $t$ carries holds on
every run of $t$.

\noindent\textbf{Proof Sketch.}
We first show that $R_\pi \subseteq D_\pi$. By (i) a run reaches its exit carrying
the path it took and not the path anyone asked about. Intuitively, an input that
slips out of $R_\pi$ into a sibling path arrives at that sibling's exit with a
record that is not $\pi$, and the assertion waiting there refutes the query. Such
an input cannot be present in a region the verifier discharged, and by (iv) a
certificate reports a region it did discharge and not one it declined.

We then show that every run of $t$ lies in $A_t$ and satisfies the assertions. By
(ii) and (iii), each coordinate the region constrains is given a value inside the
corresponding set of $A_t$ or checked to hold one before the call, so a run of
$t$ cannot leave the part the verifier was asked about;
Section~\ref{sec:oracle} proves every emitted assertion over $R_\pi$, which
contains $A_t$. Composing the two parts, every run of $t$ is one of the
executions the verifier quantified over, and the assertions hold on it.

\noindent\textbf{Corollary.} None of this depends on $\Pi_G(u)$ being complete.
A path never enumerated, one dropped when the target budget ran out, and a
decision no longer recorded because its call site was withdrawn all cost
strength. The first two leave fewer siblings, so the candidate arrives wider
and may take more rounds to certify; the third leaves a coarser class of runs
for an assertion to hold over. Intuitively, each takes information away from
the construction and none away from the check, and it is the check the
certificate rests on. Soundness therefore rests on $\Delta_G(u)$ being
instrumented completely, which is local and can be inspected point by point,
and not on the enumeration being complete, which is global.

Four limits belong with the proposition rather than with the threats to
validity. A witness names quantities of the model, and where one of them has no
value a test can write, the preamble of Section~\ref{sec:emit} cannot establish
it and the check placed before the call fails aloud; that gap is closed by a
test that stops, and not by an assumption. Where an unsettable quantity can
change the recorded path, the region leaves it free, the certification of
Section~\ref{sec:cert} is refuted or left undecided, and no region beyond the
witness is emitted. A path that exists in the model and not in the contract is
set aside in Section~\ref{sec:enum}, and certification could not have caught
it, there being no execution for a refutation to exhibit. And the verifier is
taken to be free of defects, as Section~\ref{sec:task-def} states.

\subsection{Verified Assertion Synthesis}\label{sec:oracle}

Section~\ref{sec:region} leaves a set of inputs that all follow one path, and
nothing yet says what those inputs do. This stage receives the path $\pi$, the
certified region $R_\pi$, the quantities the contract declares, and the values
the witness of Section~\ref{sec:enum} carried at its exit. It returns a set of
\emph{verified assertions}, statements the verifier has proved to hold on every
input $R_\pi$ admits. What may be stated is
decided before anything is proved, and what decides it is the set of quantities
a caller can observe.

\subsubsection{What Can Be Asserted}
An \emph{observable} of a unit is a quantity whose identity the fixture and the
path fix and whose value the emitted test can read once the call has finished.
The execution model settles the list: the balance and the stored values of each
account the call touched, the data it returned, whether it succeeded, and the
entries it logged. Nothing else outlives the call. Candidates are therefore
enumerated from what the contract declares, and not drawn from a catalogue of
properties. Where they are read is settled as well, since the path fixes the
exit the call reaches and leaves no observation point to choose.

The statements come in three layers that differ in what they cost. The first is
carried by the path itself. Which exit the call reaches is fixed by $\pi$, and
so is any reason or log entry the recorded decisions determine on their own, so
those hold on every input of $R_\pi$ with no query put at all. A test states
them as an expected revert, an expected reason, or an expected log entry; where
a path returns and leaves nothing else to name, the test states that the call
succeeded, which the framework would otherwise leave implicit. What
a withdrawn call site does inside itself is not among them, since two runs of
one class may fail different checks or log different entries there. Neither is
the content of a reason or an entry, which moves with the input. Those go to
the layers below.

The second layer says which observables the call left alone. Each is placed in
one of three states over the part in hand: unchanged on every input it admits,
changed on every one, or settled by neither answer. The vocabulary is two
comparisons, that an observable equals the value it held before the call or
that it differs from it, which is what a specification language states as a
frame condition. Each observable carries two candidates, one for each
comparison: a refuted equality shows only that some input changes the
observable, and the other candidate is what settles whether every input does.
Neither answer can be read off the text of the path, since a write need not
change the value it writes to, and a value can change inside a call the path
makes. Reads are not among the three states: a caller cannot see a read once
the call is over, so a read noted during verification has nowhere to go in a
test.

The third layer says what the changed ones became. Its statements are order
relations between an observable and the value it held before the call, and
bounds on that value or on the difference between the two. They apply to an
observable whose type orders its values, and whose order the unit itself uses.
Addresses and booleans are integers underneath, and a statement about where one
address falls relative to another holds without saying anything. The third layer
is therefore offered only where the text of the unit compares the quantity or
does arithmetic on it. The two comparisons of the second layer apply to every
observable, and for one whose type carries no equality of its own a test
compares the bytes. Those two quantities,
the final value and the change, are the only subjects. A wider vocabulary would
turn into a catalogue of shapes, and which shapes belong in it is a question
about a domain, not about the contract at hand.

\subsubsection{Candidates and Dominance}
The forms of the last two layers carry an order of strength. Writing
$\mathit{post}$ and $\mathit{pre}$ for the value of an observable after and
before the call, they are
(1)~$\mathit{post} = \mathit{pre}$;
(2)~$\mathit{post} \neq \mathit{pre}$;
(3)~$\mathit{post} \geq \mathit{pre}$;
(4)~$\mathit{post} > \mathit{pre}$;
(5)~$a \leq \mathit{post} \leq b$;
and (6)~$a \leq \mathit{post} - \mathit{pre} \leq b$.
The first two say whether the observable moved and the next two which way, each
of those two standing for its mirror as well. The last two say by how much, the
difference in (6) taken in whichever direction (3) settled. Form (1) implies
both directions of (3), and (4) implies (3) together with (2). The two bounds
neither imply nor follow from the order relations. A bound can be wide where an
order relation is exact, and narrow where none holds at all. Some forms are
therefore incomparable, and for each observable \tool{VeriPUT} keeps every form
the verifier discharges that no other discharged form implies.

No form equates an observable with a term outright. The forms above are proposed
without reading what the path computes, each being a fixed shape filled from a
fixed vocabulary. An order relation or a bound holds for many terms at once, so
an enumeration to the declared depth finds one wherever one exists. An equality
holds for a single term and fails for the terms beside it, so proposing one
means searching that space until it lands exactly. That is a different
construction, and we leave it to later work.

A form may also be written with a condition in front of it, and checked only
where that condition holds. This is for the inputs a region cannot separate.
Where the runs on which an observable moves are not a product of ranges, no
division of the region describes them. Two things bound the condition.
(i)~It draws on what a term draws on, so a condition the path computes and the
test cannot evaluate is not available. (ii)~It is not invented. Within one path
every recorded decision is fixed, so the decisions left open are those the
granularity of Section~\ref{sec:budget} does not record, and those, rewritten
over the quantities the test holds, are the conditions on offer. A withdrawn
call site returns here what it gave up at the path. Both sides are put to the
verifier separately, and a side left unproved is dropped.

The ends of (5) and (6) are written as terms and not as numbers. A term is built by a grammar of sums, differences, products and
division by a literal, to a declared depth. What it draws on is what the test
has in hand where the assertion stands: the values it rendered for the
coordinates, the values it read before the call, the constants the contract
declares, and the literals the text of the unit contains. Gathering those is a
reading of declarations and of literals. No value is substituted for another,
no flow is followed, and which terms hold is left to the verifier; the
generator never reads the expression a path computes. The depth is therefore
what keeps an emitted assertion from becoming a transcription of the
implementation, and it decides how many candidates an observable receives.

Each term must also be one the test can evaluate. The query that discharges an
assertion discharges at the same time that its terms are defined on every input
of the part, so an overflow in the oracle cannot be read as a fault in the
contract.

Terms keep two quantities apart that numbers would tie together. A bound written
as a number has to hold at the worst input its part admits, so the wider that
part the weaker the bound. What a wider region buys in coverage it would give
back in strength. A bound built from the quantities of the path does not move
when the part does, which leaves the queries spent and the tests emitted as the
only price of strength.

\subsubsection{What the Verifier Is Asked}
Each candidate is decided on its own, and the verdicts of Section~\ref{sec:enum}
apply here unchanged: discharged, refuted with an input, or declined. Candidates
do not depend on one another, so an implementation may put several to the
verifier at once, in the shape Section~\ref{sec:candidate} already used, but
nothing below rests on their being asked together. The order changes how many
are asked and not which are provable. A discharged form settles every form it
implies, and the second layer decides for each observable whether the third has
anything left to say. For a candidate $a$ and a
region $A$, what is asked is
\begin{equation}
\mathsf{Proved}(A, a) \;\equiv\;
\mathsf{Discharged}\bigl(\mathsf{assume}(A),\ \mathsf{assert}(a)\bigr),
\label{eq:oracle-query}
\end{equation}
where $\mathsf{assume}(A)$ is the same text the emitted test writes to
constrain its arguments. The verifier is asked under the same constraint on the
inputs that the emitted test renders, and not under a restatement of it. The
path plays no part in the query: $A$ lies inside $R_\pi$, which
Section~\ref{sec:cert} certified, so every run under that assumption follows
$\pi$ already.

A declined candidate is dropped and recorded with the reason it received. A
refuted one is not, where the input the verifier returned can divide the region.
That input is the split point, and on the side of it holding $x_\pi$ the
candidate is asked again. The parts cover exactly what they divided, so nothing
is given up in coverage, and the price is one further emitted test each.
The division keeps a statement that a single region would have lost, and the
budget on tests for one path is where it stops.

A path therefore yields a family of parts, each with the assertions proved over
it, drawn from $R_\pi$ by Section~\ref{sec:emit}:
\begin{equation}
R_\pi \;=\; A_1 \cup \cdots \cup A_k \cup E_\pi ,
\label{eq:parts}
\end{equation}
where each $A_i$ is the set of inputs one emitted test admits and $E_\pi$
collects the inputs left unrendered. Two things divide $R_\pi$: a candidate
refuted on one side of an input the verifier returned, and punctures on a
coordinate too dense for a rendering to skip. Every $A_i$ lies inside $R_\pi$,
so what was proved over the whole of it holds on all of them and is not asked
again, while what a division saved holds on the part that saved it. $E_\pi$
holds what the budget on tests for one path leaves, and is reported and not
absorbed.

Where no candidate at all is discharged over $R_\pi$, no parameterized test is
emitted and the concrete replay test of Section~\ref{sec:enum} stands in its
place. It is not left without an oracle. At a single input every observable has
one value, and the refutation that produced that input reported them, so the
test states those values at no further query. It is the weakest oracle the
construction emits and the strongest available at one point, and it is why every
test carries an assertion the verifier established.

\subsubsection{What Each Test Carries}
Listing~\ref{lst:skeleton} gives the shape every emitted test takes. The
contract is held in a field (Line~2) and deployed by the fixture the run settled
on (Line~4), which every test repeats. The name records the unit and the path
(Lines~1 and~6), and the parameters are the coordinates the rendering left more
than one value (Line~6). Each of those receives its range on entry, and a
coordinate the refinement pinned is written as the value it holds (Line~7). The
preamble then establishes the state and context coordinates (Line~8). The values
the statements below compare against are read before the call (Line~9). Where
the framework needs a statement of the first layer ahead of the call, it goes
there (Line~10). The call follows (Line~11). Statements proved over the whole
part come after it (Line~12), and those a condition separates carry that
condition (Line~13). A test for a reverting path carries the first layer alone,
since a revert undoes what the other two would have compared, and one that fell
back to a single input carries the values the refutation reported. Line~7
assigns and never filters, and Line~13 holds the only conditions a test
contains.

\begin{lstlisting}[language=Solidity, numbers=left, numberstyle=\scriptsize,
xleftmargin=1.5em, basicstyle=\ttfamily\scriptsize, breaklines=true,
columns=fullflexible, keepspaces=true, float=tbp,
caption={The shape of an emitted test. Slots in braces are filled per path;
those that a path leaves nothing for are dropped.},
label={lst:skeleton}]
contract {UNIT}_{PATH}_Test is Test {
    {CONTRACT} c;

    function setUp() public { {FIXTURE} }

    function test_{UNIT}_{PATH}({PARAMS}) public {
        {RANGES}      // one per coordinate: bound(p, L, U), or a fixed value
        {PREAMBLE}    // establish the state and context coordinates
        {PRE}         // read every observable the statements below compare
        {EXIT}        // first layer, where it must precede the call
        c.{UNIT}({ARGS});
        {FLAT}        // statements proved over the whole part
        if ({COND}) { {THEN} } else { {ELSE} }   // statements a condition separates
    }
}
\end{lstlisting}

Take the withdrawal of Listing~\ref{lst:motivation} on the path where a discount
applies and the fee stays under the cap. From the first layer, the call returns
instead of reverting and logs one withdrawal. From the second, the owner, the
account the fee is paid to and the caller's discount are as they were, while the
caller's deposit and the balance of the account the fee reaches are not. From
the third, the deposit falls by exactly the amount passed in, and the fee lies
between that amount taken at the lowest rate the discount range leaves in force
and the same amount taken at the highest.

$M_2$ charges the rate on the caller's whole deposit rather than on the amount
withdrawn, and refutes the upper of those two bounds where the deposit is large
enough against the amount, which is a proper part of the region. $M_1$
multiplies every fee by a hundred, so the subtraction computing the amount
returned underflows and the call reverts, which refutes the statement of the
first layer on every input the region admits.

A developer writing this test by hand might state the amount returned exactly,
as Listing~\ref{lst:motivation-put} states two bounds on the fee.
\tool{VeriPUT} states a bound instead, for the reason given in
Section~\ref{sec:oracle}. Weakening an exact value into a bound has a price, and
this is where it is paid.

An assertion may relate one coordinate to another. The restriction of
Section~\ref{sec:candidate} is on the region, which a framework has to render by
assigning each coordinate a value of its own. A statement checked once the call
is over is under no such constraint, and the last of the three above names both
the amount passed in and the amount returned. Each assertion a test carries was
proved under the assumption that test itself writes, which is what
Section~\ref{sec:sound-emit} takes from this stage.

\subsection{Implementation}\label{sec:implementation}

\begin{answerbox}
\noindent\textbf{Summary.}
\end{answerbox}

\section{Evaluation}\label{sec:eval}

\begin{rqbox}
\begin{enumerate}
  \item[\textbf{RQ1}] (Performance)
  \item[\textbf{RQ2}] (Failure study)
  \item[\textbf{RQ3}] (Ablation study)
  \item[\textbf{RQ4}] (Robustness study)
\end{enumerate}
\end{rqbox}

\subsection{Benchmarks and Environment Setup}\label{sec:setup}

\subsubsection{Dataset Selection Criteria}

\subsubsection{<Subject categories with counts>}

\subsubsection{Pre-processing}\label{sec:preprocess}

\subsubsection{Environment}

\subsection{RQ1: Performance}
\begin{answerbox}\textbf{Answer to RQ1.}\end{answerbox}

\subsection{RQ2: Failure Study}
\begin{answerbox}\textbf{Answer to RQ2.}\end{answerbox}

\subsection{RQ3: Ablation Study}
\begin{answerbox}\textbf{Answer to RQ3.}\end{answerbox}

\subsection{RQ4: Robustness across <backends>}
\begin{answerbox}\textbf{Answer to RQ4.}\end{answerbox}

\subsection{Case Study}\label{sec:case_study}

\subsubsection{Case Study 1: <the motivation example>}

\subsubsection{Case Study 2: <the complementary situation>}
\begin{answerbox}\textbf{Summary.}\end{answerbox}

\section{Threats to Validity}\label{sec:threats}

We summarize the main threats to validity of this study.

\subsection{Internal Validity}
\textbf{<Threat>.}

\subsection{External Validity}
\textbf{<Threat>.}

\subsection{Construct Validity}
\textbf{<Threat>.}

\subsection{Statistical Conclusion Validity}
\textbf{<Threat>.}

\section{Related Work}\label{sec:related}

A generated unit test needs a criterion that settles which tests the suite
contains, an input part, and a part that decides whether a run was correct. We
review coverage-driven test generation for smart contracts
(Section~\ref{sec:rw-cov}), the synthesis of parameterized unit tests
(Section~\ref{sec:rw-put}), and the inference of invariants and properties of
contracts (Section~\ref{sec:rw-inv}).

\subsection{Coverage-driven test generation for smart contracts}\label{sec:rw-cov}

For smart contracts, \tool{SolAR}~\cite{solar_sbst}, developed from the earlier
\tool{AGSolT} framework~\cite{agsolt}, and
\tool{SynTest-Solidity}~\cite{syntest_solidity} are search-based tools that
evolve test cases toward higher branch coverage. \tool{SolAR} emits a sequence
of transactions with concrete arguments, while \tool{SynTest-Solidity} executes
each candidate on a local chain through \tool{Truffle} and records the values
observed in that execution as the assertions of the test it emits.
\tool{SolTG}~\cite{SolTG-1} instead encodes the contract's executions as
constrained Horn clauses, enumerates the path constraints of that encoding one
function at a time and for call sequences of growing length, solves each with
an SMT solver, and writes the arguments it obtains into \tool{Foundry} test
functions, stopping once every branch is covered. \tool{CC-SolBMC}~\cite{CCSolBMC2024}
works below the branch: at each atomic condition it plants a trap assertion,
one that fails exactly when the condition takes a given value, and takes the
counterexample a bounded model checker returns as the test input for that
condition. \tool{TestSmart}~\cite{TestSmart2021} draws its
test data from an off-the-shelf symbolic execution engine and contributes a
mutation-based reduction of that engine's output, and
\tool{Griffin}~\cite{Griffin2024} takes a target line as part of its input and
returns the data that reaches it.

Beyond smart contracts, trap assertions are the usual device for driving a
verifier toward a coverage target, and three lines of work bear on ours. Wei et
al.~\cite{wei2025esbmc} plant one per branch in the GOTO program of a C program
and compute coverage from the verifier's answers, but do not address turning the
counterexamples into tests. Seeding Contradiction~\cite{SeedingContradiction2025}
does collect them, into a regression suite for an Eiffel program, planting one
trap assertion per branch in the Boogie translation and requiring the program to
have been proved correct beforehand so that the only failures are the ones it
planted. \tool{DTest}~\cite{DTest2023} plants one at a time in the Boogie
translation of a Dafny program and turns each counterexample into a unit test,
filling the arguments the counterexample leaves unconstrained with default
values and taking the oracle from the postcondition of the method under test
where the developer wrote one.

\tool{VeriPUT} differs from these tools in three ways. First, its coverage
target is a complete decision sequence rather than a branch, an atomic
condition, or a basic block. A target of the latter kind is reached in the
middle of a run, so the counterexample witnesses that the target is reachable
and says nothing about how the run ends. Criteria that keep a path as the
target but bound its length fall short on the same count, and on a second.
Among them are $k$-limiting path coverage and the prime path coverage recently
added to a production compiler as a
measurement~\cite{AmmannOffutt,primepath_gcc}. A prime path is a maximal simple
path, so it need not reach an exit, and a single run covers many prime paths at
once, so the targets do not divide the inputs into classes that a test can
assume. Second, these tools fix every argument of the test they emit, so one
path is exercised at one point of the inputs that follow it, and where an
argument is left unconstrained a default value stands in for the range it could
have taken; \tool{VeriPUT} emits a region of inputs and proves that every input
in it follows the same path. Third, none of them synthesizes an assertion
proved over a range of inputs. An emitted test either reports a fault only when
the fault stops the run, transcribes the values one execution produced, or
inherits an assertion the developer already wrote, as in \tool{DTest}, whose
tests for a method without a postcondition are left with none. We compare
against a generator when it (i) takes the contract source as its only input and
(ii) produces a suite of its own tests rather than a witness for a target that
the user supplies. \tool{Griffin} fails both. \tool{TestSmart} meets both, but
its implementation is not available, and its own contribution is a
mutation-based reduction of a suite that a separate symbolic execution engine
produces, so the generator a comparison would measure is that engine.

\subsection{Synthesizing parameterized unit tests}\label{sec:rw-put}

To the best of our knowledge, no prior work synthesizes parameterized unit
tests for smart contracts, so we review this line beyond the contract domain.
Tillmann and Schulte~\cite{Tillmann2005} introduced the parameterized unit
test, in which the developer writes both the assumption that constrains the
arguments and the assertion that judges the run, and later work derives one or
both of these from tests that already exist~\cite{Danglot2019}. The assumption
is the easier of the two. Thummalapenta et al.~\cite{Thummalapenta2011} leave it
to the developer and supply a methodology for writing it,
\tool{JARVIS}~\cite{Peleg2018Jarvis} fits a template drawn from a library of
abstract domains to a group of tests that repeat one scenario, and
\tool{PROZE}~\cite{Proze2024} takes the arguments a method is observed to
receive at runtime, while \tool{Teralizer}~\cite{Glock2025} extracts the path
condition of one execution symbolically, which is exact but requires the method
under test to be a pure function. The assertion is harder, because deriving one
presupposes a notion of what the code ought to do, and most of this line
therefore reuses the assertion the developer already wrote, keeping it when it
survives the widened inputs~\cite{Peleg2018Jarvis,Proze2024}. Fraser and
Zeller~\cite{FraserZeller2011} are the exception: they generalize the
postcondition of a concrete test into candidates and retain those that mutation
analysis fails to refute.

\tool{VeriPUT} differs from these tools in three ways. First, it too widens a
concrete input into a region, but the input it widens is a counterexample it
obtained in order to reach a path that the coverage criterion named, whereas
these methods widen a test the developer had already written, so the paths their
output reaches are the paths that test reached. Second, they establish the
widened assumption by running it. Sampling and mutation analysis can refute a
candidate but cannot establish one, so what these methods report is that the
inputs they tried behaved as expected; what \tool{VeriPUT} reports is that a
verifier has proved the property for every input the assumption admits. Third,
exact extraction, as in \tool{Teralizer}, buys its exactness with a purity
requirement that the entry points of interest here do not meet, since they read
and write persistent state by construction.

\subsection{Invariant and property inference for smart contracts}\label{sec:rw-inv}

For smart contracts, \tool{InvCon}~\cite{InvCon2022} mines the transaction
history of a deployed contract for likely invariants, relations over its state
variables and function arguments reported at the level of the contract and of
individual functions, and \tool{InvCon+}~\cite{InvConPlus2024} adds a static
phase so that the relations it reports are proved rather than likely.
\tool{SmartOracle}~\cite{SmartOracle2025} mines the same histories for
finer-grained relations over the values a function reads and writes, organized
per branch, and checks them against later transactions.
\tool{Trace2Inv}~\cite{Trace2Inv2024} instead fixes the shape of the invariant
in advance, instantiating a set of templates from history and deploying the
result as a guard that rejects offending transactions. Where no history is
available the evidence is the code itself: \tool{SmartInv}~\cite{SmartInv2024}
learns invariants from the code together with its comments in order to locate
vulnerabilities, and \tool{PropertyGPT}~\cite{PropertyGPT2025} retrieves
human-written properties to produce formal ones for a prover to discharge.
Beyond smart contracts, \tool{Daikon}~\cite{Daikon2007} reports the properties
that hold across the runs of an existing suite, and supplies the categories,
such as a value left unchanged or one variable bounded by another, that this
line reuses.

None of these tools generates tests, so none can serve as a baseline: they emit
no inputs, they carry no coverage criterion, and what they produce is an
annotation for a prover, a guard deployed beside the contract, or a report of a
suspected vulnerability, none of which a test framework can run. Two further
differences bear on the assertions \tool{VeriPUT} synthesizes. First, they work
in a different mode: most are on-chain analyses whose evidence is the
transaction history of a contract already deployed, whereas \tool{VeriPUT} runs
off-chain before deployment and reads only the source. Second, the statements
sit at different levels of abstraction. An invariant has to survive every entry
point and every order in which calls may arrive; what \tool{VeriPUT} attaches to
a test is a property of one path through one entry point. A path-scoped
statement need only hold on its path and may fail elsewhere, while a
contract-wide invariant holds across all paths and does not by itself
characterise the behaviour of any one of them.

\section{Conclusion}\label{sec:conclusion}

\section{Ethics, Responsible Disclosure, and Data Availability}

\bibliographystyle{ACM-Reference-Format}
\bibliography{ref}

\end{document}